% chap08_numpy.tex - Numerical computing with NumPy
% ============================================================================

\chapter{NumPy \& Numba: High Speed Numerical Computing}
\label{chap:numpy}

Python lists are flexible but slow for numerical work. NumPy (Numerical Python) provides arrays---efficient containers for numerical data---and fast operations on them. NumPy is the foundation of scientific Python; virtually every scientific library builds on it. You can install it from your terminal, preferably with a virtual environment active, via:

\begin{lstlisting}
pip install numpy
\end{lstlisting}

% ----------------------------------------------------------------------------
\section{Why NumPy?}
\label{sec:why-numpy}

Compare adding two lists of numbers:

\begin{lstlisting}
# Pure Python - slow
a = [1, 2, 3, 4, 5]
b = [10, 20, 30, 40, 50]
result = [x + y for x, y in zip(a, b)]

# NumPy - fast
import numpy as np
a = np.array([1, 2, 3, 4, 5])
b = np.array([10, 20, 30, 40, 50])
result = a + b  # array([11, 22, 33, 44, 55])
\end{lstlisting}

NumPy is faster because:
\begin{itemize}
    \item Arrays store data in contiguous memory
    \item Operations are implemented in optimized C code
    \item Operations work on entire arrays at once (``vectorization'')
\end{itemize}

For large datasets, NumPy can be 100x faster than pure Python.

% ----------------------------------------------------------------------------
\section{Creating Arrays}
\label{sec:creating-arrays}

Import NumPy with the conventional alias:

\begin{lstlisting}
import numpy as np
\end{lstlisting}

\subsection{From Lists}

\begin{lstlisting}
# 1D array
a = np.array([1, 2, 3, 4, 5])

# 2D array (matrix)
b = np.array([[1, 2, 3],
              [4, 5, 6]])

# Note: all elements must be the same type
c = np.array([1, 2, 3.0])  # all become floats
print(c.dtype)  # float64
\end{lstlisting}

\subsection{With Functions}

\begin{lstlisting}
# Zeros and ones
np.zeros(5)           # array([0., 0., 0., 0., 0.])
np.ones((2, 3))       # 2x3 array of ones
np.full((2, 2), 7)    # 2x2 array filled with 7

# Ranges
np.arange(0, 10, 2)   # array([0, 2, 4, 6, 8])
np.linspace(0, 1, 5)  # array([0., 0.25, 0.5, 0.75, 1.])

# Identity matrix
np.eye(3)             # 3x3 identity matrix

# Random numbers
np.random.rand(3)           # 3 uniform random numbers [0, 1)
np.random.randn(3)          # 3 standard normal random numbers
np.random.randint(0, 10, 5) # 5 random integers from 0-9
\end{lstlisting}

\begin{tip}
\py{np.linspace(start, stop, num)} creates exactly \py{num} evenly spaced points including both endpoints. \py{np.arange(start, stop, step_size)} creates values from \py{start} to \py{end-step_size}, with values differenting by \py{step_size}. Which is better really depends on usecase and needs.
\end{tip}

% ----------------------------------------------------------------------------
\section{Array Properties}
\label{sec:array-properties}

\begin{lstlisting}
a = np.array([[1, 2, 3],
              [4, 5, 6]])

a.shape      # (2, 3) - 2 rows, 3 columns
a.ndim       # 2 - number of dimensions
a.size       # 6 - total number of elements
a.dtype      # dtype('int64') - data type
\end{lstlisting}

Common data types:
\begin{itemize}
    \item \py{int32}, \py{int64}: Integers
    \item \py{float32}, \py{float64}: Floating point (64 is default)
    \item \py{bool}: Boolean
    \item \py{complex128}: Complex numbers
\end{itemize}

Specify type when creating:
\begin{lstlisting}
a = np.array([1, 2, 3], dtype=np.float32)
\end{lstlisting}

\begin{tip}
The numbers (32, 64) indicate how many \textbf{bits} (1s and 0s) are used to store each value. A \texttt{float64} uses twice the memory of a \texttt{float32}, allowing for much higher precision and larger ranges. For example, an \textbf{int8} can only hold numbers from -128 to 127. If you try to store 500, it simply won't fit! Meanwhile, int16 can store -32,768 to 32,767. While 64-bit floats are generally a safe default, the curious can find the technical details of how these are stored by searching for \href{https://en.wikipedia.org/wiki/Floating-point_arithmetic}{``Floating-point arithmetic.''}, or following that link to the wikipedia page on the topic.
\end{tip}

% ----------------------------------------------------------------------------
\section{Indexing and Slicing}
\label{sec:numpy-indexing}

NumPy indexing extends Python's list indexing:

\subsection{1D Arrays}

\begin{lstlisting}
a = np.array([10, 20, 30, 40, 50])

a[0]        # 10 (first element)
a[-1]       # 50 (last element)
a[1:4]      # array([20, 30, 40])
a[::2]      # array([10, 30, 50]) - every other
\end{lstlisting}

\subsection{2D Arrays}

\begin{lstlisting}
b = np.array([[1, 2, 3],
              [4, 5, 6],
              [7, 8, 9]])

b[0, 0]      # 1 (row 0, col 0)
b[1, 2]      # 6 (row 1, col 2)
b[0]         # array([1, 2, 3]) - first row
b[:, 0]      # array([1, 4, 7]) - first column
b[0:2, 1:3]  # 2x2 subarray: [[2, 3], [5, 6]]
\end{lstlisting}

\subsection{Boolean Indexing}

Select elements that meet a condition:

\begin{lstlisting}
a = np.array([1, 5, 3, 8, 2, 9])

a[a > 4]           # array([5, 8, 9])
a[a % 2 == 0]      # array([8, 2]) - even numbers

# Modify based on condition
a[a > 5] = 0       # set values > 5 to 0
print(a)           # array([1, 5, 3, 0, 2, 0])
\end{lstlisting}

\begin{labexample}
Filtering spectroscopy data:

\begin{lstlisting}
wavelengths = np.linspace(400, 700, 100)
absorbances = np.random.rand(100) * 0.5

# Find wavelengths where absorbance > 0.3
mask = absorbances > 0.3
high_abs_wavelengths = wavelengths[mask]
print(f"High absorbance at {len(high_abs_wavelengths)} wavelengths")
\end{lstlisting}
\end{labexample}

% ----------------------------------------------------------------------------
\section{Array Operations}
\label{sec:array-operations}

\subsection{Element-wise Operations}

Operations apply to every element:

\begin{lstlisting}
a = np.array([1, 2, 3, 4])
b = np.array([10, 20, 30, 40])

a + b       # array([11, 22, 33, 44])
a - b       # array([-9, -18, -27, -36])
a * b       # array([10, 40, 90, 160])
a / b       # array([0.1, 0.1, 0.1, 0.1])
a ** 2      # array([1, 4, 9, 16])
np.sqrt(a)  # array([1., 1.41..., 1.73..., 2.])
\end{lstlisting}

\subsection{Mathematical Functions}

Vectorized functions are orders of magnitude faster for large volumes of high dimensional data than functions applied in loops would be. Scripts that can take hours using loops might taken seconds with vectorized functions. NumPy provides vectorized versions of many common math functions, examples include: 

\begin{lstlisting}
x = np.array([0, np.pi/4, np.pi/2, np.pi])

np.sin(x)      # sine
np.cos(x)      # cosine
np.exp(x)      # e^x
np.log(x)      # natural log (watch out for log(0)!)
np.log10(x)    # base-10 log
np.abs(x)      # absolute value
\end{lstlisting}

\begin{tip}
It is a good idea to search up "numpy documentation" to have a look at what other functions and operations are avalaible through numpy. In your programming journey, getting comfortable with finding and making use of documentation is among the most important skills you will have. It is well beyond the scope of this book to give an exhaustive insight into all NumPy has to offer.
\end{tip}

\subsection{Aggregation Functions}

Summarize array data:

\begin{lstlisting}
a = np.array([1, 5, 3, 8, 2])

np.sum(a)      # 19
np.mean(a)     # 3.8
np.std(a)      # standard deviation
np.var(a)      # variance
np.min(a)      # 1
np.max(a)      # 8
np.argmin(a)   # 0 (index of minimum)
np.argmax(a)   # 3 (index of maximum)
\end{lstlisting}

For 2D arrays, specify axis:

\begin{lstlisting}
b = np.array([[1, 2, 3],
              [4, 5, 6]])

np.sum(b)           # 21 (all elements)
np.sum(b, axis=0)   # array([5, 7, 9]) - sum each column
np.sum(b, axis=1)   # array([6, 15]) - sum each row
\end{lstlisting}

% ----------------------------------------------------------------------------
\section{Broadcasting}
\label{sec:broadcasting}

NumPy can operate on arrays of different shapes through \emph{broadcasting}:

\begin{lstlisting}
a = np.array([1, 2, 3])
b = 10

a + b    # array([11, 12, 13]) - scalar broadcast to all elements
a * b    # array([10, 20, 30])
\end{lstlisting}

More complex broadcasting:

\begin{lstlisting}
# 2D array + 1D array
matrix = np.array([[1, 2, 3],
                   [4, 5, 6]])
row = np.array([10, 20, 30])

matrix + row
# array([[11, 22, 33],
#        [14, 25, 36]])
\end{lstlisting}

The 1D array is ``broadcast'' across each row of the 2D array.

\begin{labexample}
Baseline subtraction:

\begin{lstlisting}
# Spectrum with baseline
wavelengths = np.linspace(400, 700, 100)
raw_spectrum = np.random.rand(100) + 0.2  # signal + offset

# Subtract baseline (minimum value)
baseline = np.min(raw_spectrum)
corrected = raw_spectrum - baseline  # broadcasts scalar
\end{lstlisting}
\end{labexample}

% ----------------------------------------------------------------------------
\section{Reshaping Arrays}
\label{sec:reshaping}

\begin{lstlisting}
a = np.arange(12)  # array([0, 1, 2, ..., 11])

# Reshape to 3x4
b = a.reshape(3, 4)
# array([[ 0,  1,  2,  3],
#        [ 4,  5,  6,  7],
#        [ 8,  9, 10, 11]])

# Reshape to 2x6
c = a.reshape(2, 6)

# Flatten back to 1D
d = b.flatten()  # or b.ravel()

# Transpose (swap rows and columns)
e = b.T
# array([[ 0,  4,  8],
#        [ 1,  5,  9],
#        [ 2,  6, 10],
#        [ 3,  7, 11]])
\end{lstlisting}

\begin{warning}
\py{reshape()} requires the total number of elements to stay the same. You can't reshape a 12-element array into 3x5 (15 elements).
\end{warning}

% ----------------------------------------------------------------------------
\section{Stacking and Splitting}
\label{sec:stacking}

\begin{lstlisting}
a = np.array([1, 2, 3])
b = np.array([4, 5, 6])

# Stack vertically (row by row)
np.vstack([a, b])
# array([[1, 2, 3],
#        [4, 5, 6]])

# Stack horizontally (column by column)
np.hstack([a, b])
# array([1, 2, 3, 4, 5, 6])

# Concatenate along an axis
np.concatenate([a, b])  # same as hstack for 1D
\end{lstlisting}

Splitting:

\begin{lstlisting}
c = np.arange(12).reshape(3, 4)

# Split into 3 parts along rows
np.vsplit(c, 3)  # list of three 1x4 arrays

# Split into 2 parts along columns
np.hsplit(c, 2)  # list of two 3x2 arrays
\end{lstlisting}

% ----------------------------------------------------------------------------
\section{Linear Algebra}
\label{sec:linear-algebra}

NumPy provides linear algebra operations:

\begin{lstlisting}
a = np.array([[1, 2],
              [3, 4]])
b = np.array([[5, 6],
              [7, 8]])

# Matrix multiplication
np.dot(a, b)     # or a @ b
# array([[19, 22],
#        [43, 50]])

# Transpose
a.T

# Determinant
np.linalg.det(a)  # -2.0

# Inverse
np.linalg.inv(a)

# Eigenvalues and eigenvectors
eigenvalues, eigenvectors = np.linalg.eig(a)

# Solve linear system Ax = b
A = np.array([[3, 1], [1, 2]])
b = np.array([9, 8])
x = np.linalg.solve(A, b)  # x such that A @ x = b
\end{lstlisting}

\begin{labexample}
Linear regression using linear algebra:

\begin{lstlisting}
# Data: concentration vs absorbance
conc = np.array([0.1, 0.2, 0.3, 0.4, 0.5])
absorbance = np.array([0.12, 0.25, 0.35, 0.48, 0.61])

# Set up design matrix for y = mx + c
# [1 x1]       [c]     [y1]
# [1 x2]   @   [m]  =  [y2]
# ...                   ...
A = np.column_stack([np.ones_like(conc), conc])
b = absorbance

# Solve least squares: (A^T A) x = A^T b
coeffs = np.linalg.lstsq(A, b, rcond=None)[0]
intercept, slope = coeffs

print(f"Beer's Law: A = {slope:.3f} * C + {intercept:.3f}")
print(f"Molar absorptivity (slope): {slope:.3f} L/(mol*cm)")
\end{lstlisting}
\end{labexample}

% ----------------------------------------------------------------------------
\section{Saving and Loading Data}
\label{sec:numpy-io}

\begin{lstlisting}
data = np.random.rand(100, 5)

# NumPy binary format (fast, preserves precision)
np.save("data.npy", data)
loaded = np.load("data.npy")

# Text format (human-readable)
np.savetxt("data.csv", data, delimiter=",")
loaded = np.loadtxt("data.csv", delimiter=",")

# Save multiple arrays
np.savez("multiple.npz", wavelengths=wl, absorbances=ab)
loaded = np.load("multiple.npz")
print(loaded["wavelengths"])
\end{lstlisting}
% ----------------------------------------------------------------------------
\section{Einstein Summation}
\label{subsec:einsum}

One of NumPy's most powerful (and often overlooked) tools is \py{np.einsum}. This function uses Einstein summation notation to represent many common array operations using a short string of letters. It is an advanced skill to developed, and is mostly added here to let you know it exists.

The notation identifies the axes of the input arrays with letters. If a letter appears in the input but not the output, that axis is summed over.

\begin{lstlisting}
a = np.array([1, 2, 3])
b = np.array([10, 20, 30])
A = np.arange(6).reshape(2, 3)

# Dot product: sum(a_i * b_i) -> output is scalar
np.einsum('i,i', a, b)       # 140

# Matrix-vector multiplication: sum(A_ij * a_j) -> output is vector (i)
np.einsum('ij,j->i', A, a)   # array([ 8, 26])

# Sum all elements in a matrix
np.einsum('ij->', A)         # 15

# Transpose
np.einsum('ij->ji', A)       # Same as A.T
\end{lstlisting}

\begin{tip}
While \py{np.einsum} can be intimidating at first, it is often more memory-efficient than standard operations because it avoids creating intermediate temporary arrays. If you find yourself nesting multiple \py{np.reshape}, \py{np.transpose}, and \py{np.dot} calls, \py{einsum} is likely the cleaner solution and worth the effort learning. For a deep dive into the logic of indices, I highly recommend exploring the official \href{https://numpy.org/doc/stable/reference/generated/numpy.einsum.html}{NumPy documentation} on the topic.
\end{tip}
% ----------------------------------------------------------------------------
\section{Examples of Common Patterns}
\label{sec:numpy-patterns}

\subsection{Normalize Data}

\begin{lstlisting}
data = np.array([10, 20, 30, 40, 50])

# Scale to 0-1 (MinMax Scaling)
normalized = (data - data.min()) / (data.max() - data.min())

# Z-score normalization
z_scores = (data - data.mean()) / data.std()
\end{lstlisting}

\subsection{Moving Average}

\begin{lstlisting}
def moving_average(data, window):
    return np.convolve(data, np.ones(window)/window, mode='valid')

smoothed = moving_average(noisy_data, window=5)
\end{lstlisting}

\subsection{Finding Peaks}

\begin{lstlisting}
from scipy.signal import find_peaks # scipy is covered more in a later section

data = np.array([1, 3, 7, 1, 2, 6, 4, 2])
peaks, _ = find_peaks(data, height=5)  # indices where peaks > 5
print(f"Peaks at indices: {peaks}")
\end{lstlisting}

% ----------------------------------------------------------------------------
\newpage
\section{When Vectorization Isn't An Option: Numba}
\label{sec:numba}

NumPy's vectorization is powerful, but some algorithms resist vectorization, or perhaps you are unsure of how to vectorize your logic. Perhaps you need nested loops that depend on previous results, or complex branching logic that doesn't map to array operations. In these cases \py{Numba} may offer you some performance gains. It compiles your Python code to machine code on-the-fly, giving you C-like speed while writing Python. The first step of a loop is therefore slow, but subsequent steps are blazing fast, so for long-running loops it can really add up. You can install it with:

\begin{lstlisting}[style=shellstyle]
pip install numba
\end{lstlisting}

\subsection{How Numba Works}
\label{sec:Numba}
Numba uses JIT (Just-In-Time) compilation. You decorate (recall Section~\ref{sec:decorators-dataclasses}) a function with \py{@numba.jit}, and Numba converts it to compiled code the first time you call it. Subsequent calls run at compiled speed. The catch: Numba only works with NumPy arrays and basic Python operations. Complex features like classes or dictionary look-ups won't compile.

\subsection{A Comprehensive Example: Molecular Reactions}

Let's simulate a system of molecules diffusing and reacting upon collision. In this model, we assume that any unreacted molecule will react if it gets close enough to another unreacted molecule during a timestep. Importantly, because we evaluate the entire system before updating any states, molecules can react in clusters. This means if one molecule is near several others simultaneously, they will all react together. This requires checking the distances between all pairs and updating the reactions synchronously at the end of the timestep. This is logic that can be vectorized, but it still serves as an excellent proving ground for alternative optimization approaches. 

Over the new few pages you will see four approaches, each required to show the complete picture. To begin, you'll need to import the following:

\begin{lstlisting}
import math
import random
import time
import numpy as np
from numba import jit
\end{lstlisting}

\begin{note}
The specific chemical simulation explored here serves as a vehicle to demonstrate performance differences. If the physical chemistry of diffusing molecules doesn't align with your exact interests, you need not worry, it is not imperative to understand every line of the loop logic. However, if you do wish to follow along, I have included detailed comments to help guide you. 

The key lessons to take away from the following examples are:

\begin{itemize}
    \item \py{NumPy} is not automatically faster just by importing it; unvectorized array access can actually slow you down.
    \item Vectorization can be unintuitive and filled with contextual pitfalls, making it hard to offer universal "rules of thumb" for.
    \item Numba provides an excellent middle ground, by allowing you to write straightforward and readable loops while achieving lightning-fast compiled execution speeds.
\end{itemize}
\end{note}

\newpage 
\noindent \textbf{Example 1: Pure Python with Lists}
\begin{lstlisting}
def react_pure_python(n_molecules, time_steps, react_distance):

	# Randomly initialise positions
    pos = [[random.random() * 100, random.random() * 100] for _ in range(n_molecules)]
	
	# Mark all molecules as un-reacted
    reacted = [False] * n_molecules

	
    for t in range(time_steps):
        # Diffusion, everything steps in a random direction a small amount
        random_steps = [random.gauss(0, 0.1) for _ in range(n_molecules * 2)]
        for i in range(n_molecules):
            pos[i][0] += random_steps[i * 2]
            pos[i][1] += random_steps[i * 2 + 1]

        # Track the reactions for this timestep
        new_reacted = [False] * n_molecules
        for i in range(n_molecules):
            if reacted[i]: continue
            for j in range(i + 1, n_molecules):
                if reacted[j]: continue
                
                dx = pos[i][0] - pos[j][0]
                dy = pos[i][1] - pos[j][1]
                dist = math.sqrt(dx*dx + dy*dy)
                
                if dist < react_distance:
                    new_reacted[i] = True
                    new_reacted[j] = True

        # Synchronous update: Apply all reactions at end of timestep
        for k in range(n_molecules):
            if new_reacted[k]:
                reacted[k] = True

    return sum(reacted)

start = time.time()
result = react_pure_python(n_molecules=500, time_steps=10_000, react_distance=2.0)
print(f"Pure Python (lists): {time.time() - start:.2f} seconds ({result} reacted)")
\end{lstlisting}

\noindent On a typical machine: \textbf{6-10 seconds}. Baseline performance with interpreted Python.

\newpage 
\noindent \textbf{Example 2: NumPy (Unvectorized --- Same loop but with \py{Numpy} arrays)}

\begin{lstlisting}
def react_numpy_unvectorized(n_molecules, time_steps, react_distance):
    """Unvectorized: NumPy arrays, but same loop structure as pure Python."""
    pos = np.random.rand(n_molecules, 2) * 100
    reacted = np.zeros(n_molecules, dtype=np.bool_)

    for t in range(time_steps):
        # Diffusion (batch generation)
        pos += np.random.randn(n_molecules, 2) * 0.1

        # Track reactions for this timestep
        new_reacted = np.zeros(n_molecules, dtype=np.bool_)
        for i in range(n_molecules):
            if reacted[i]: continue
            for j in range(i + 1, n_molecules):
                if reacted[j]: continue
                
                dx = pos[i, 0] - pos[j, 0]
                dy = pos[i, 1] - pos[j, 1]
                dist = np.sqrt(dx*dx + dy*dy)
                
                if dist < react_distance:
                    new_reacted[i] = True
                    new_reacted[j] = True

        # Synchronous update
        reacted |= new_reacted

    return np.sum(reacted)

start = time.time()
result = react_numpy_unvectorized(n_molecules=500, time_steps=10_000, react_distance=2.0)
print(f"NumPy (unvectorized): {time.time() - start:.2f} seconds ({result} reacted)")
\end{lstlisting}

\noindent On a typical machine: \textbf{15-29 seconds}. Slower than pure Python! NumPy without vectorization can actually be slower due to array access overhead in tight loops. Clearly we can't simply assume using NumPy makes our code faster.

\newpage 
\noindent \textbf{Example 3: NumPy (Vectorized --- Better NumPy Usage)}

\begin{lstlisting}
def react_numpy_vectorized(n_molecules, time_steps, react_distance):
    """Vectorized: compute all pairwise distances at once."""
    pos = np.random.rand(n_molecules, 2) * 100
    reacted = np.zeros(n_molecules, dtype=np.bool_)

    for t in range(time_steps):
        # Diffusion (vectorized)
        pos += np.random.randn(n_molecules, 2) * 0.1

        # Reactions (vectorized distance matrix)
        if not np.all(reacted):
            active_indices = np.where(~reacted)[0]
            active_pos = pos[active_indices]
            
            # Compute distances for active molecules only
            diff = active_pos[:, np.newaxis, :] - active_pos[np.newaxis, :, :]
            distances = np.sqrt(np.sum(diff**2, axis=2))

            # Diagonal matrix trick to ignore self-interaction (distance 0)
            np.fill_diagonal(distances, np.inf) 

            # Find which rows in the SUBSET have a neighbor within distance
            reacting_in_subset = np.any(distances < react_distance, axis=1)

            # Map those subset results back to the original 'reacted' array
            # This is inherently synchronous
            reacted[active_indices[reacting_in_subset]] = True

    return np.sum(reacted)

start = time.time()
result = react_numpy_vectorized(n_molecules=500, time_steps=10_000, react_distance=2.0)
print(f"NumPy (vectorized): {time.time() - start:.2f} seconds ({result} reacted)")

\end{lstlisting}

\noindent On a typical machine: \textbf{$\approx$1 second}. There are two main tricks here to make it all work: the first is to prevent self-interaction, and the second is to shrink the distance matrix by dropping molecules that are done. Without applying both tricks, the answers will either be wrong or take much longer than the other methods ($>$ 1 minute). While this example is still somewhat simple, vectorizing more complex processes that we can easily describe using built-in Python logic can become quite a headache. So, what if we want speed, but don't want to undertake an entire mathematics degree to get it? Well, we can explore example four:

\newpage 
\noindent \textbf{Example 4: NumPy + Numba (Compiled Unvectorized)}

\begin{lstlisting}
@jit(nopython=True)
def react_numba(n_molecules, time_steps, react_distance):
    """Compiled: Numba accelerates the unvectorized loop logic."""
    pos = np.random.rand(n_molecules, 2) * 100
    reacted = np.zeros(n_molecules, dtype=np.bool_)

    # Pre-allocate outside the loop for maximum C-level performance
    new_reacted = np.zeros(n_molecules, dtype=np.bool_)

    for t in range(time_steps):
        # Diffusion
        pos += np.random.randn(n_molecules, 2) * 0.1

        # Reset the temporary array for this timestep
        new_reacted.fill(False)
        
        for i in range(n_molecules):
            if reacted[i]: continue
            for j in range(i + 1, n_molecules):
                if reacted[j]: continue
                
                dx = pos[i, 0] - pos[j, 0]
                dy = pos[i, 1] - pos[j, 1]
                dist = np.sqrt(dx*dx + dy*dy)
                
                if dist < react_distance:
                    new_reacted[i] = True
                    new_reacted[j] = True

        # Synchronous update
        for k in range(n_molecules):
            if new_reacted[k]:
                reacted[k] = True

    return np.sum(reacted)

# Warmup run to hide JIT compilation time in benchmarks (uncomment the line below for a no-overhead estimate)
#react_numba(n_molecules=10, time_steps=1, react_distance=2.0)

start = time.time()
result = react_numba(n_molecules=500, time_steps=10_000, react_distance=2.0)
print(f"Numba (compiled, warmed-up): {time.time() - start:.2f} seconds ({result} reacted)")
\end{lstlisting}

\noindent On a typical machine: \textbf{1-2 seconds} with the cold-start compilation overhead, or \textbf{$<$0.45 seconds} if warmed up (see the commented-out warmup call above). The first invocation includes the cost of compiling to machine code; subsequent calls run at compiled speed. For a one-off script, vectorized NumPy edges ahead because there is no compilation penalty. For repeated calls---parameter sweeps, Monte Carlo studies, interactive analysis---Numba's compiled speed dominates once the initial cost is paid. Both are roughly \textbf{10x faster} than pure Python, but they achieve it differently: NumPy eliminates loops through array operations, Numba keeps the loops but compiles them to machine code.

\begin{note}
\textbf{The Complete Picture (real timings, 500 molecules, 10,000 timesteps):}
\begin{itemize}
    \item \textbf{Pure Python (lists):} $\sim$10 seconds --- baseline interpreted Python
    \item \textbf{NumPy unvectorized:} $\sim$17 seconds --- \emph{slower} than pure Python due to array access overhead in tight loops
    \item \textbf{NumPy vectorized:} $\sim$1 second --- fast, when done carefully (see Section~\ref{sec:ai-usage} for what happens when it isn't)
    \item \textbf{Numba compiled:} $\sim$1.7 seconds cold-start, $\sim$\textbf{0.4 seconds} warmed up
\end{itemize}

\textbf{Key lessons:}
\begin{enumerate}
    \item \textbf{Vectorization wins when done thoughtfully.} But ``thoughtfully'' is doing a lot of work in that sentence. Getting the vectorized example right required excluding self-interactions, subsetting to active molecules, and understanding the memory implications. The naive version was the \emph{slowest} of all four approaches.
    \item \textbf{Numba is the pragmatic choice} when vectorization is too complex or when you need repeated calls. Write straightforward loops, add a decorator, and get compiled speed.
    \item Using NumPy without vectorization is \textbf{worse than not using NumPy at all.} Array element access in Python loops carries overhead that raw Python lists avoid.
    \item \textbf{Understanding your problem determines the right tool.} There is no universal hierarchy---the best approach depends on problem structure, call frequency, and how much time you are willing to invest in getting the implementation right.
\end{enumerate}
\end{note}

\begin{tip}
\textbf{Best Practices for Numba Compatibility}

Numba compiles Python to machine code, but not all Python features are supported. Keep your code simple to avoid compilation errors:

\begin{itemize}
    \item \textbf{Use basic loops and conditionals} --- \py{for}, \py{while}, \py{if}/\py{else} work great. Avoid Python's functional tools (\py{map}, \py{filter}, \py{lambda}).

    \item \textbf{Pre-allocate arrays, don't append} --- Numba doesn't support dynamic lists. Use \py{np.empty()} and track indices (as in our kinetics example), rather than appending.

    \item \textbf{Stick to basic NumPy operations} --- Use array indexing, arithmetic, and functions like \py{np.sqrt()}, \py{np.sum()}. Avoid fancy indexing or complicated reshape logic.

    \item \textbf{Avoid strings and dictionaries} --- These don't compile. No string operations, no dictionary lookups, no print statements in compiled code.

    \item \textbf{Use \py{nopython=True}} --- Force Numba to compile or raise an error immediately, rather than silently falling back to slow Python mode: \py{@jit(nopython=True)}.

    \item \textbf{Test your algorithm without Numba first} --- Get the logic right in pure Python, then add the decorator. Python errors are clearer than Numba compilation errors.
\end{itemize}

When in doubt, keep the code as simple as possible. Numba works best on straightforward, loop-heavy code using NumPy arrays and basic math.
\end{tip}

\subsection{When to Use Numba}

Use Numba when:
\begin{itemize}
    \item You have loops over NumPy arrays that can't be vectorized
    \item The logic is complex (nested loops, conditionals, sequential dependencies)
    \item Speed matters (simulation, expensive analysis)
\end{itemize}

Don't use Numba when:
\begin{itemize}
    \item Your code already uses NumPy vectorization well
    \item The bottleneck is I/O or calling external libraries
    \item You need Python-only features (classes, dictionaries, dynamic types)
\end{itemize}

\begin{tip}
Always profile before optimizing. Use the \py{time} module or \py{timeit} to measure. As we saw above, ``using NumPy'' does not automatically mean ``fast''---the unvectorized NumPy example was slower than pure Python. Similarly, getting vectorization right took careful thought about subsetting and self-interactions. Measure first, understand the problem, optimize second.
\end{tip}

\begin{warning}
Because Numba resulted in such a dramatic speedup for our collision simulation, it is tempting to conclude that Numba is a consistent upgrade. This is absolutely false. There are many scenarios where pure, vectorized NumPy is significantly faster than anything you can write in Numba:

\begin{itemize}
    \item \textbf{Reinventing the Wheel (Sorting and Searching):} It is incredibly common for beginners to write custom loops to sort an array, find unique values, or search for a specific threshold. However, built-in functions like \py{np.sort()}, \py{np.unique()}, and \py{np.searchsorted()} are backed by highly optimized C algorithms. Compiling your own sorting loop with Numba will almost certainly be slower than just calling NumPy's native version.
    \item \textbf{Complex Aggregations:} If you find yourself writing a loop to bin data into a histogram, compute a cumulative sum, or find the median across a specific axis, check the documentation first. Functions like \py{np.histogram()}, \py{np.cumsum()}, and \py{np.median()} already operate at peak hardware efficiency.
    \item \textbf{Heavy Linear Algebra:} Operations like matrix multiplication (\py{A @ B}), matrix inversions, and eigenvalue decompositions are handed off by NumPy to heavily optimized, multi-threaded external libraries (like OpenBLAS or Intel MKL). A Numba-compiled \py{for} loop cannot compete with decades of hardware-specific tuning.
\end{itemize}

\textbf{The golden rule:} If NumPy already has a dedicated function for exactly what you want to do, use it! Only reach for Numba when you have custom, complex, or branching iterative logic (like tracking states over time) that resists standard vectorization.
\end{warning}

% ----------------------------------------------------------------------------
\section{Chapter Summary}
\label{sec:numpy-summary}

NumPy is essential for scientific computing:

\begin{itemize}
    \item \textbf{Arrays:} \py{np.array()}, \py{np.zeros()}, \py{np.linspace()}
    \item \textbf{Properties:} \py{.shape}, \py{.dtype}, \py{.size}
    \item \textbf{Indexing:} \py{a[0]}, \py{a[1:5]}, \py{a[a > 0]}
    \item \textbf{Operations:} Element-wise (\py{+}, \py{*}), functions (\py{np.sin}, \py{np.exp})
    \item \textbf{Aggregation:} \py{np.sum()}, \py{np.mean()}, \py{np.std()}
    \item \textbf{Linear algebra:} \py{@} or \py{np.dot()}, \py{np.linalg.solve()}
    \item \textbf{Reshaping:} \py{.reshape()}, \py{.T}, \py{np.vstack()}
    \item \textbf{Einsum Notation:} for more complex operations
    \item \textbf{Numba:} \py{@jit(nopython=True)} compiles loop-heavy code to machine code---ideal when vectorization is too complex or when amortising compilation cost across repeated calls
\end{itemize}

\begin{ponder}
\begin{enumerate}
    \item Why is \py{a * b} element-wise multiplication for arrays but \py{a @ b} is matrix multiplication?
    \item How would you select all values in an array that are between 5 and 10?
    \item What does \py{axis=0} vs \py{axis=1} mean for a 2D array?
    \item Write code to normalize an array so its values sum to 1.
\end{enumerate}
\end{ponder}
